\begin{proposicion}[Concentración de los observables]\label{prop:concentracion-meso}
Para una familia fija de patrones y tipos de fibras de pesos acotados, fijemos $b_1>0$ y $\chi>0$.
Existe $c_{\rm pat}>0$ tal que, para cada suelo $K$, en las historias de entrada regular, salvo
probabilidad $\le C e^{-c_{\rm pat}K}$, cada observable alcanzado antes de la primera salida o en la
transición que
llega a ella difiere de su valor en $M$ en a lo sumo
\[
 C_H\{A(\delta+1/n)+\chi+1/K\}+b_1 .
\]
Para las fibras se compara cada línea distinguida todavía disponible con su propia fibra en $M$. La
tasa depende de las constantes de la clase exterior, de la familia de observables, de $b_1$ y de
$\chi$.
\end{proposicion}
\begin{proof}
Restamos a cada incremento su esperanza predecible y congelamos el proceso a partir del tablero de
primera salida. Para un patrón sin línea distinguida, el salto centrado en el paso desde el tamaño
$u$ es a lo sumo $2C_H/u$. La desigualdad maximal de Azuma \eqref{macro:max-mart}, con
$\sum_{u>k}u^{-2}\le1/k$, acota por $2\exp(-c_H b_1^2k)$ la probabilidad de que la martingala alcance
en módulo $b_1$ desde $M$ hasta $k$. El
Lema~\ref{lema:presupuesto-meso} y la suma de $u^{-2}$ acotan la deriva acumulada en todos los
prefijos a la vez.

Fijada la historia de entrada y una de sus líneas distinguidas, usemos el núcleo auxiliar que en cada
padre regular divide el núcleo original restringido a los hijos que mantienen disponible la línea
distinguida por su probabilidad de supervivencia $y_u$. Este núcleo sólo condiciona el paso siguiente.
En cualquier prefijo que llega con la línea distinguida disponible al tamaño $k$, incluido el que
termina en la primera salida, el peso bajo $\nu$ es el peso auxiliar multiplicado por
\[
 \prod_{u=k+1}^{M}y_u
 =\prod_{u=k+1}^{M}\big(1-1/u+O(u^{-2})\big)
 =(k/M)\exp(O(1/k)).
\]
Para la fila el producto es exactamente $k/M$. Para la columna el error es uniforme aun cuando los
factores dependan de los tableros del prefijo. La misma identidad de pesos vale en un instante de
parada porque se aplica por separado a cada prefijo terminal. Bajo el núcleo auxiliar restamos la
deriva condicional del Lema~\ref{lema:incrementos-meso} y congelamos al llegar al tablero de primera
salida. La
misma desigualdad maximal de Azuma vale desde cualquier historia de entrada regular.

En la banda $m\le k<2m$, con $m=2^jK$, paremos en la primera observación de la banda en que la martingala
de esta línea distinguida tenga módulo $\ge b_1$, si existe. Al transferir los prefijos de parada a
$\nu$, el factor cuesta a lo sumo $(2m/M)e^{C/m}$. La unión condicional sobre las $M$ líneas
distinguidas de cada familia y los tipos fijos de pesos, y después la suma sobre las bandas, cuestan a
lo sumo $\sum_m C m e^{-c_H b_1^2 m}\le C'e^{-c' b_1^2K}$. Mientras la línea distinguida siga disponible,
la resta de deriva en estas trayectorias es exactamente la del núcleo auxiliar; su módulo acumulado
queda controlado por el presupuesto en la trayectoria original. El fallo del presupuesto se cuenta una
sola vez bajo $\nu$, mediante el Lema~\ref{lema:presupuesto-meso}, después de integrar sobre la
historia de entrada. Así no aparece un factor $M$ en la cola de ese fallo ni un factor $n$ en la cola
de las fibras.
\end{proof}

\subsection{Reconstrucción de las condiciones de regularidad}\label{meso:regularidad}
Escribimos $W_0=\nu_1/(u_Du_A)=b_0r_Rr_C$, que es el valor en $x=0$ del peso $W_x$ de la
Proposición~\ref{macro:trayectoria}, $q_{ij}=Y_{ij}W_{0,ij}/k$ y
$V_{ij}=\nu_{1,ij}/k$, evaluados en las coordenadas originales de la casilla. Las cotas $c_\nu,C_\nu$
para la densidad, la vacancia mínima $\kappa_\nu$ y la constante de Lipschitz $L_\nu$ son las
constantes fijas del perfil. Consideramos una partición fija de $[0,1]$ en intervalos de longitud a
lo sumo $a_{\rm mes}$. El radio $\varrho_{\rm mac}$ y el paso $a_{\rm mac}$ de la
Sección~\ref{macro} se eligen después, en el punto~5 de la Sección~\ref{meso:constantes}, en función
de esta partición, de las tolerancias que siguen y de $\delta$ (Lema~\ref{meso:entrada}). Para filas
y columnas controlamos sus proporciones restantes en cada intervalo con objetivo la longitud del
intervalo. Fijamos una tolerancia $\varrho_{\rm mes}>0$ y tolerancias individuales cuya suma es a lo
sumo $\varrho_{\rm mes}$ por familia, por ejemplo $\varrho_{\rm mes}|I|$ para $I$ en esta partición.
Las fibras de fila y columna tienen pesos $W_0$, objetivo $1$ y tolerancia a lo sumo
$\varrho_{\rm mes}$ cada una.

Para describir los otros patrones escribamos $T=\{1,2\}\times\{1,2\}$. Por cada $S\subseteq T$ tomamos
el patrón de dos filas y dos columnas que requiere una casilla disponible en las posiciones de $S$. Su
peso es $\prod_{e\in S}W_{0,e}\prod_{e\in T\setminus S}\nu_{1,e}$. Los dieciséis tienen un mismo
objetivo,
\[
 \int_{[0,1]^4}\nu_1(r,c)\nu_1(r,c')\nu_1(r',c)\nu_1(r',c')\,dr\,dr'\,dc\,dc' .
\]
La normalización inyectiva es $(k)_2^2$. Usar $k^4$ y agregar las asignaciones que repiten fila o
columna modifica cada patrón en $O(1/k)$, porque hay $O(k^3)$ asignaciones adicionales con peso
acotado. Con $E=q-V$, como en el Lema~\ref{macro:brecha}, la identidad
\[
 \operatorname{tr}\big[(EE^{T})^2\big]
 =\sum_{i,i',j,j'} E_{ij}E_{i'j}E_{i'j'}E_{ij'}
\]
produce los dieciséis términos así extendidos con signos $(-1)^{4-|S|}$. Los dieciséis objetivos se
cancelan por $(1-1)^4=0$. Se escogen las tolerancias positivas por separado de modo que su suma, más
$O(1/K_0)$, sea menor que la cuarta potencia de la cota espectral requerida. Entonces $\|q-V\|_{\rm
op}\le\operatorname{tr}[(EE^T)^2]^{1/4}$ es menor que dicha cota. La expansión incluye
$S=\varnothing$, y permite diagonales repetidas en todos los términos.
